\section{The exact minimum for the projected graph normalization}
\label{sec:w-exact-minimum}

The preceding reduction permits a complete, finite computation for the auxiliary normalization $w=e_W[\Gamma_1]$. This section concerns the complex product $E^4$ with its fixed product polarization. Its computational certificate uses exact integral cohomology, and the geometric completeness of the catalog is supplied by W01--W03.

\begin{theorem}[W04]
Let $E=\C/\Z[\omega]$, let $A=E^4$, let $\theta$ be the product principal polarization, and use the Weil action $\eta(k)=\operatorname{diag}(k,k,\bar k,\bar k)$. For the normalization $w=e_W[\Gamma_1]$ of P06 one has
\[
\boxed{c_{\C,\theta}(w)=24.}
\]
The lower bound is a computer-assisted finite proof: W01--W03 reduce every possible certificate of cost below twenty-four to the abelian-surface catalog specified below, and the accompanying self-contained exact-integer program exhausts that catalog. The upper bound is the explicit six-surface identity below and is independent of the negative search.
\end{theorem}

\begin{proof}
Write $q=3w$. The exact integral Weil lattice computed in V07 shows that $w$ has denominator three. Thus an equality $[Z^+]-[Z^-]=Nw$ with integral cycles requires $3\mid N$. The complex calibration of P08 gives
\[
N=|z(Nw)|\leq \frac{\deg_\theta Z^++\deg_\theta Z^-}{8}.
\]
For cost below twenty-four, $N$ must consequently equal three. Since $\theta^2q=0$, both effective numerators have the same degree, say $D$. Surface degrees for this product polarization are even. The previously proved lower bound $16\leq c_{\C,\theta}(w)$, or the period classification below, leaves only
\[
D\in\{16,18,20,22\}.
\]

W02 replaces every hypothetical numerator by an effective integral sum of translated abelian surfaces with the identical cohomological class and degree. Indeed each numerator has nonzero total Weil coordinate, since $|z(Z)|\leq22/8<3$ and their difference has coordinate three. Every zero-Weil component therefore has degree at most fourteen. W02 applies to those components and to all nonzero-Weil components of degree at most twenty-two. Translations do not affect cohomology, so they introduce no parameter in the finite comparison.

We first specify the catalog of nonzero-Weil components. By W03 every such component is a graph or transposed graph of an Eisenstein matrix $h$, with
\[
d=\Nm(\det h)\in\{1,3,4\},\qquad
T=\sum_{i,j}\Nm(h_{ij})\leq10-d,\qquad
\deg_\theta\Gamma_h=2(1+d+T).
\]
All entries are therefore in the explicitly bounded list
\[
\{a+b\omega:\ -3\leq a,b\leq3,\quad a^2-ab+b^2\leq9\}.
\]
Both orientations are included for $d=3,4$. For $d=1$ a transposed graph is already the graph of the inverse Eisenstein matrix and is included automatically.

Every nonzero-Weil component has degree at least eight. A numerator of degree at most twenty-two therefore has at most two such components, counting multiplicities. If it has two, both are automorphism graphs: a component with $d=3$ or $4$ has degree at least sixteen or eighteen, respectively. In a two-component numerator each component has degree at most fourteen. If there is one component, it may have any of the three determinant norms above.

For a graph the period is $z(\Gamma_h)=\overline{\det h}$, and for a transposed graph it is its conjugate. Thus the total period of a numerator belongs to the finite set of Eisenstein integers of norm $0,1,3,$ or $4$: it is either a single element of norm $1,3,$ or $4$, or a sum of two units. Requiring $z^+-z^-=3$ leaves exactly
\[
(z^+,z^-)=
(1,-2),\quad(2,-1),\quad
(1-\omega,-2-\omega),\quad(2+\omega,-1+\omega).
\]
Applying $[-1]$ to the first elliptic coordinate sends $q$ to $-q$ and preserves degree. Applying it to both numerators and then interchanging them replaces the first pair by the second. It is therefore enough to retain
\[
\mathcal P=\{2,1-\omega,2+\omega\},\qquad
\mathcal M=\mathcal P-3.
\]
The program also checks this finite list directly in Eisenstein arithmetic. A positive numerator has either two automorphism components or one component with determinant norm three or four. A negative numerator has either two automorphism components or one component whose period lies in $\mathcal M$; the latter includes period $-1$ automorphism graphs. This accounts for every retained case.

We next specify all zero-Weil components of degree at most fourteen. For every one of the six coordinate pairs, enumerate its graphs with
\[
\sum_{i,j}\Nm(h_{ij})+\Nm(\det h)\leq6,
\]
retaining exactly those with zero Weil coordinate. Each entry lies in
\[
\{a+b\omega:\ -2\leq a,b\leq2,\quad a^2-ab+b^2\leq5\};
\]
the norm six does not occur, so no entry allowed by the preceding budget is lost. W03 proves that the only remaining abelian surfaces are the degree-fourteen exceptions
\[
E_i\times\{(a t,b t)\in E_j\times E_k:t\in E\}\times\{0\}_{\ell},
\qquad\{\Nm(a),\Nm(b)\}=\{3,4\}.
\]
Their equations, used to compute the classes, are
\[
x_\ell=0,\qquad b x_j-a x_k=0.
\]
The coefficients $a,b$ are coprime, so the second equation has connected kernel and introduces no unrecorded covering multiplicity. All coordinate choices and unit associates are included. Deduplication uses the entire integral cohomology class. The resulting catalog has the following numbers of classes:
\[
\begin{array}{c|rrrrrrr}
\text{degree}&2&4&6&8&10&12&14\\\hline
\text{zero-Weil abelian classes}&6&72&288&180&1440&864&1440.
\end{array}
\]
In particular, the last entry includes the $144$ classes without a degree-one coordinate-pair projection. A graph-only catalog would omit them.

All effective zero-Weil sums of degree at most fourteen are generated by an exact finite recurrence. If $\mathcal G_d$ is the set of generator classes of degree $2d$, process $d=1,\ldots,7$ in turn. For each target degree $2n$, adjoin every sum of an already generated class of degree $2(n-kd)$ and a multiset of $k$ elements of $\mathcal G_d$, for $0\leq k\leq\lfloor n/d\rfloor$. The empty class starts degree zero. This permits every nonnegative multiplicity and imposes no bound on the number of generators other than degree. The class counts at degrees $0,2,\ldots,14$ are
\[
1,\quad6,\quad93,\quad776,\quad6066,\quad48564,\quad322050,\quad2162952.
\]

For each $D\in\{16,18,20,22\}$, the program adds every possible zero-Weil remainder to each retained positive main part and records the resulting classes of degree exactly $D$. It then enumerates every retained negative main part with every zero-Weil remainder of the complementary degree and asks whether its class plus $q$ appears in the positive set. The positive main part always has degree at least sixteen, whereas a negative period-$-1$ graph may have degree eight. Thus remainders through degree fourteen, including the exceptional surfaces, are necessary and sufficient. The exact counts and outcomes are
\[
\begin{array}{c|r|r|c}
D&\text{distinct positive classes}&\text{negative candidates tested}&\text{equality found}\\\hline
16&750&193344&\text{none}\\
18&14364&1544688&\text{none}\\
20&184902&11659752&\text{none}\\
22&1898952&79234032&\text{none}.
\end{array}
\]
The negative count includes duplicates, which affect only runtime. No equality is discarded by such duplicates.

For completeness, the exact class representation and its comparison rule are as follows. Use the ordered integral one-form basis
\[
(u_1,v_1,u_2,v_2,u_3,v_3,u_4,v_4)
\]
and the seventy increasing exterior monomials in degree four. Multiplication by $a+b\omega$ pulls back the pair $(u_i,v_i)$ to
\[
(a u_i-b v_i,\ b u_i+(a-b)v_i).
\]
The class of a graph is the exterior product of its four integral defining one-forms; the exceptional class is obtained from the two displayed homomorphism equations. These formulae compute all seventy coefficients, not only the Weil projection. For
\[
P=D_{13}(1),\quad Q=D_{13}(\omega),\quad
U=D_{24}(1),\quad V=D_{24}(\omega),
\]
the target is computed as
\[
q=PU-PV-QU-2QV.
\]
The program independently verifies
\[
18q=(9+S)(2-R-R^2)[\Gamma_1],\qquad q^2=18,
\qquad\theta^2q=0,
\]
so that its target agrees exactly with the normalization $q=3w$ in P06.

For efficient storage the vector $(c_0,\ldots,c_{69})$ is encoded by the integer $\sum_i c_i1001^i$. The largest absolute coefficient of any retained generator is seven; the largest absolute target coefficient is two. A cycle of degree at most twenty-two has at most eleven irreducible components, with multiplicities. Hence every coefficient of the difference of two compared cycles minus $q$ has absolute value at most
\[
2\cdot11\cdot7+2=156<1001.
\]
If the encoded difference vanishes, its lowest coefficient is a multiple of $1001$ of absolute value less than $1001$, so it vanishes. Division by $1001$ and induction prove that all seventy coefficients vanish. This establishes exact injectivity for every comparison, including the target subtraction and all effective sums. The program uses arbitrary-precision integers and no approximate comparisons. Its ordinary set membership compares full integers; hash collisions do not identify unequal integers.

The self-contained certificate \texttt{research/w\_exact24\_certificate.py} regenerates every catalog from the displayed finite lists, checks the coefficient bound before the main comparison, and reproduces the table. It requires only Python's standard library, and no external data or saved binary catalog. Its complete output is retained as \texttt{research/w\_exact24\_certificate.out}. Together with the proved exhaustive reductions above, the negative comparison excludes every certificate of cost below twenty-four.

Finally put $Z_{a,b}=\{(x,y,ax,by):(x,y)\in E^2\}$. The following integral cohomological identity supplies the upper bound:
\begin{align*}
3w={}&[Z_{1,-\omega}]+[Z_{1,-\omega^2}]+[Z_{\omega,-1}]\\
&-[Z_{-\omega,-\omega}]-[Z_{-\omega^2,-1}]-[Z_{-\omega^2,-\omega^2}].
\end{align*}
Each of these six surfaces has degree eight. Its positive and negative numerators therefore have degree twenty-four, and its multiplier is three. The certificate verifies this identity in all seventy integral coefficients before performing the negative search. Thus $c_{\C,\theta}(w)\leq24$, completing the proof.
\end{proof}

\paragraph{Scope of specialization.}
W04 is the exact complex result for this particular normalization and polarization. The degree-twenty-four identity specializes wherever the product, its CM endomorphisms, and the relevant cycle classes have good reduction. The eventual-constancy theorem B03 then gives $c_v(w)=24$ outside a finite set of good places, including along both the ordinary and supersingular subsequences. W04 does not claim that every exceptional good prime has the same minimum; no new classification of exceptional supersingular fibers is used. The all-good-prime results for $t_*$ and $t_0$ remain the separate theorems P13 and V06.
